\documentclass[12pt,a5paper]{article}

% Russian language support
\usepackage[T2A]{fontenc}
\usepackage[utf8]{inputenc}
\usepackage[russian]{babel}

% Page layout
\usepackage[a5paper,margin=1.5cm]{geometry}
\usepackage{setspace}
\setstretch{1.06}

% Graphics and figures
\usepackage{graphicx}
\usepackage{float}

% Tables
\usepackage{longtable}
\usepackage{booktabs}
\usepackage{array}
\usepackage{tabularx}

% Mathematics
\usepackage{amsmath}
\usepackage{amssymb}
\usepackage{amsthm}
\newcommand{\Kcal}{\mathcal{K}}

% Hyperlinks
\usepackage{hyperref}
\hypersetup{
    colorlinks=true,
    linkcolor=black,
    citecolor=black,
    urlcolor=black
}

% Headers and footers
\usepackage{fancyhdr}
\pagestyle{fancy}
\fancyhf{}
\renewcommand{\headrulewidth}{0pt}
\fancyfoot[C]{\thepage}

% Section formatting
\usepackage{titlesec}
\titleformat{\section}{\centering\normalsize\bfseries}{\thesection}{1em}{}
\titleformat{\subsection}{\normalsize\bfseries}{\thesubsection}{1em}{}

% Reduce spacing
\titlespacing*{\section}{0pt}{8pt}{4pt}
\titlespacing*{\subsection}{0pt}{6pt}{3pt}

% Paragraph settings
\setlength{\parindent}{1.25cm}
\setlength{\parskip}{0pt}

\begin{document}
\sloppy

% Suppress page number on first page
\thispagestyle{empty}

% ============================================================
% TITLE PAGE
% ============================================================
\begin{center}
\textbf{Национальный исследовательский ядерный университет <<МИФИ>>}

\vspace{1cm}

\textit{На правах рукописи}

\vspace{2cm}

\textbf{Анаэдевха Роджер Ник}

\vspace{1.5cm}

{\large \textbf{Разработка состязательно устойчивых
моделей на основе искусственного интеллекта в гибридных
системах обнаружения и предотвращения вторжений для
сетевой безопасности}}

\vspace{1.5cm}

\textbf{2.3.1.\ Системный анализ, управление и обработка\\
информации, статистика\\
(технические науки)}

\vspace{2cm}

\textbf{АВТОРЕФЕРАТ}

\vspace{0.5cm}

диссертации на соискание учёной степени\\
кандидата технических наук

\vfill

Москва -- 2026
\end{center}

\newpage

% ============================================================
% PAGE 2 — Committee and defence details
% ============================================================
\thispagestyle{empty}
\enlargethispage{1.5cm}

{\scriptsize
\begin{center}
\textbf{Работа выполнена в федеральном государственном автономном
образовательном учреждении высшего образования
<<Национальный исследовательский ядерный университет <<МИФИ>>{}>>}
\end{center}

\vspace{0.15cm}

\noindent\begin{tabular}{@{}p{2.8cm}p{7.4cm}@{}}
\textbf{Научный руководитель:} & \textbf{Трофимов Александр Геннадьевич} \newline
кандидат технических наук, доцент, доцент кафедры №\,22 <<Кибернетика>> Национального исследовательского ядерного университета <<МИФИ>> \\
\end{tabular}

\vspace{10.0cm}



\noindent Автореферат разослан \underline{\hspace{2.5cm}} 2026 года.

\vspace{0.2cm}

\noindent\begin{tabular}{@{}p{4.2cm}p{2cm}p{3.5cm}@{}}
Учёный секретарь диссертационного совета, кандидат физико-математических наук & & Климанов Сергей Геннадьевич \\
\end{tabular}
}

\newpage

% ============================================================
% MAIN CONTENT BEGINS
% ============================================================
\setcounter{page}{3}

\section*{ОБЩАЯ ХАРАКТЕРИСТИКА РАБОТЫ}
\addcontentsline{toc}{section}{Общая характеристика работы}

\subsection*{Актуальность темы исследования и степень её разработанности}

Гибридные системы обнаружения и предотвращения вторжений
(\emph{англ.}~Hybrid Intrusion Detection and Prevention Systems,
далее~--- \emph{гибридная~СОПВ}~--- система обнаружения и
предотвращения вторжений) являются основным защитным контуром
информационно-телекоммуникационных систем критического
назначения: сетей сетевого трафика, сегментов центров обработки
данных, бортовых систем беспилотных летательных аппаратов
(БПЛА), систем автономного транспорта, банковско-финансовых
платформ, медицинских информационных систем и систем
управления критическими инфраструктурами. В отличие от
традиционных коммерческих продуктов (Snort, Suricata, Zeek,
Cisco~IPS, Palo~Alto, McAfee), опирающихся на сигнатурный
анализ и фиксированные правила, в настоящей работе под
\emph{гибридной~СОПВ} понимается \emph{ИИ/МО-усиленный}
(\emph{искусственный интеллект / машинное обучение})
защитный контур, объединяющий три канонических типа
обнаружения~--- сигнатурный, поведенческий (аномальный) и
протокольно-аналитический~--- в единую многоуровневую
архитектуру с обучаемыми компонентами и интегрированными
средствами реагирования.

В качестве \emph{ключевого обучаемого компонента} гибридной
СОПВ в работе рассматривается \emph{графовая нейронная сеть}
(ГНС, \emph{англ.}~Graph Neural Network)~--- модель,
оперирующая одновременно с двумя различными входами:
числовой \emph{матрицей признаков узлов}
$\mathbf{X}\in\mathbb{R}^{n\times d}$ (где $n$~--- количество
узлов графа, $d$~--- размерность признакового вектора) и
\emph{матрицей смежности} $\mathbf{A}$, кодирующей структуру
попарных связей. ГНС формирует своё решение, распространяя
информацию из $\mathbf{X}$ \emph{вдоль} структуры $\mathbf{A}$
через итеративную агрегацию окрестностей. Принципиальная
уязвимость заключается в том, что возмущение \emph{любого}
из двух входов~--- признаков или структуры~--- может нарушить
выходное предсказание. Эта \emph{двойная поверхность атаки}
и определяет специфику задачи обеспечения устойчивости
гибридной СОПВ на основе ГНС.

Анализ инцидентов последних пяти лет в сетевой
кибербезопасности, защите бортовых систем БПЛА и автономного
транспорта показывает, что устойчивость развёрнутых
гибридных СОПВ отстаёт от их детектирующей точности на
25--65~процентных пунктов: системы, демонстрирующие 95--98\%
обнаружения на эталонных данных, теряют 20--40\% точности под
адверсариальными возмущениями входов, дрейфом распределения
трафика, действиями злонамеренных (\emph{византийских})
участников в распределённой архитектуре обучения и при
изменении профилей протоколов. К фундаментальным
неэффективностям относятся: запаздывание сигнатур на
30--90~дней относительно появления новых угроз, интенсивность
ложных срабатываний $10^{-2}$--$10^{-1}$ на один нормальный
сетевой поток и квадратичная по числу протоколов $P$ и
сессий $S$ вычислительная стоимость
$\mathcal{O}(P\cdot S)$ протокольного анализа в реальном
времени (здесь $\mathcal{O}(\cdot)$~--- общепринятое
обозначение асимптотической сложности).

В этой связи настоящая работа посвящена \emph{разработке
состязательно устойчивых моделей на основе искусственного
интеллекта в гибридных системах обнаружения и предотвращения
вторжений для сетевой безопасности}.

Единая формализация гибридной СОПВ выполнена в виде кортежа
\[
\mathcal{H} = (\mathcal{V}_t, \mathcal{E}_t, X_t, A_t,
f_\theta, g_\psi, \pi),
\]
где: $\mathcal{V}_t$~--- изменяющееся во времени множество
сущностей (узлов графа); $\mathcal{E}_t$~--- изменяющееся
во времени множество связей (рёбер); $X_t$~--- матрица
признаков узлов в момент~$t$; $A_t$~--- множество активных
сигнатур и правил; $f_\theta$~--- обучаемый детектирующий
модуль с вектором параметров~$\theta$; $g_\psi$~--- модуль
оценки уверенности с параметрами~$\psi$; $\pi$~--- политика
реагирования. Шесть сквозных свойств гибридной СОПВ~---
состязательная устойчивость (П1), робастность (П2),
эффективность (П3), распределённость (П4), человеческий
фактор (П5) и нестационарность (П6)~--- формализованы как
инвариантные условия, систематически декомпозированы на
двенадцать подзадач в виде матрицы $3\times4$
<<условия$\times$требования>> (три условия эксплуатации:
состязательная среда, нестационарность, распределённая
динамика; четыре требования: устойчивость, точность,
эффективность, приватность) и решены семью оригинальными
моделями и алгоритмами М1--М7.

Большинство существующих ГНС работают \emph{в дискретном
времени}, принудительно привязывая нерегулярные события
графа к фиксированной тактовой сетке и теряя как информацию,
так и гарантии устойчивости. Формулировки \emph{непрерывного
времени} на основе \emph{нейронных обыкновенных
дифференциальных уравнений}
(ОДУ; \emph{англ.}~Neural Ordinary Differential Equations)
гладко следуют за временем событий, охватывают множественные
временные масштабы и допускают вывод сертифицированной
устойчивости через \emph{неравенство Грёнуолла}:
$\|\mathbf{h}_1(T)-\mathbf{h}_2(T)\|\le e^{L_g T}
\|\mathbf{x}_1-\mathbf{x}_2\|$,
где $\mathbf{h}_i(T)$~--- скрытое представление при входе
$\mathbf{x}_i$ в момент~$T$, а $L_g$~--- \emph{константа
Липшица} векторного поля $f_\theta$ (равномерная верхняя
граница его коэффициента растяжения). Распределённая
среда~--- в которой несколько организаций владеют разными
фрагментами одного и того же графа и не имеют права делиться
данными~--- требует \emph{федеративного обучения} с
формальными гарантиями приватности и устойчивости к
\emph{византийским} (произвольным злонамеренным) обновлениям
от части участников. Современные модели на основе
трансформеров масштабируются квадратично по длине
последовательности $L$~--- $\mathcal{O}(L^2)$, однако
\emph{структурированные модели пространства состояний}
(\emph{англ.}~Selective State-Space Models, SSM) снижают
эту сложность до $\mathcal{O}(L\log L)$. Кроме того,
реальные графовые данные \emph{дрейфуют} после развёртывания,
и статические модели незаметно теряют точность без
сигнала об этом.

Существенный вклад в области, образующие методологический
фундамент работы, внесли отечественные и зарубежные
исследователи: в обнаружение и предотвращение вторжений и
сетевую безопасность~--- А.\,А.~Бурлаков, А.\,С.~Марков,
В.\,В.~Сачков, M.~Roesch, V.~Paxson; в графовое глубокое
обучение~--- T.\,N.~Kipf, M.~Welling, P.~Veli\v{c}kovi\'{c},
W.\,L.~Hamilton; в нейронные ОДУ и модели пространства
состояний~--- R.\,T.\,Q.~Chen, A.~Gu, T.~Dao; в адверсариальное
машинное обучение и анализ устойчивости ГНС~--- A.~Madry,
N.~Carlini, D.~Wagner, D.~Z\"ugner, S.~G\"unnemann; в
федеративное обучение с приватностью~--- B.~McMahan,
K.~Bonawitz; в теорию игр и сертификацию устойчивости~---
A.~Sinha, J.~Steinhardt, J.~Cohen. Однако \emph{ни один} из
существующих подходов не обеспечивает \emph{одновременного}
выполнения четырёх требований~--- устойчивости, точности,
эффективности и приватности~--- во всех трёх операционных
условиях гибридной СОПВ. Именно это противоречие между
многомерными требованиями и текущим состоянием алгоритмической
базы определяет актуальность настоящего исследования.

\subsection*{Объект и предмет исследования}

\textbf{Объектом исследования} являются \emph{гибридные
системы обнаружения и предотвращения вторжений (гибридная
СОПВ)}, функционирующие в состязательных, нестационарных и
распределённых средах и адаптируемые к различным предметным
областям~--- защите сетевого трафика, бортовых систем БПЛА,
бортовых систем автономного транспорта.

\textbf{Предметом исследования} являются модели, алгоритмы и
специальное математическое обеспечение, повышающие
устойчивость и эффективность гибридной СОПВ при сохранении
состязательной устойчивости, приватности распределённых
данных и калиброванности оценок неопределённости, включая
аппарат непрерывно-временной графовой динамики, селективных
моделей пространства состояний, федеративной оптимизации с
дифференциальной приватностью, оптимального транспорта,
байесовского вывода и теоретико-игровой сертификации.

\subsection*{Цель и задачи исследования}

\textbf{Целью} диссертационного исследования является
\emph{повышение устойчивости и эффективности гибридных
систем обнаружения и предотвращения вторжений (гибридная
СОПВ), функционирующих в состязательных, нестационарных и
распределённых условиях}, за счёт разработки единого
математического описания гибридной СОПВ и согласованного
семейства моделей и алгоритмов М1--М7, а также доказательства
их применимости путём инстанцирования в секторе анализа
сетевого трафика и переноса в смежный сектор защиты бортовых
систем БПЛА.

\textbf{Задачи:}

\begin{enumerate}
\item Провести систематический обзор подходов, применяемых
в современных гибридных СОПВ, и выявить их недостатки и
уязвимости в прикладных секторах.

\item Выполнить обзор методов повышения устойчивости к
состязательным атакам и обеспечения вычислительной
эффективности, включая сравнительный анализ современных
коммерческих ИИ-усиленных продуктов кибербезопасности;
сформулировать единые требования к разрабатываемым моделям
и алгоритмам.

\item Формализовать математическую модель гибридной СОПВ в
виде кортежа $\mathcal{H} = (\mathcal{V}_t, \mathcal{E}_t,
X_t, A_t, f_\theta, g_\psi, \pi)$ и шести сквозных свойств
П1--П6 в виде инвариантных условий, образующих матрицу
$3\times4$ <<условия$\times$требования>>; в явном виде
поставить многокритериальную задачу повышения устойчивости,
точности, эффективности и приватности.

\item Разработать алгоритмы и модели М1--М7 повышения
устойчивости и эффективности гибридной СОПВ в распределённых
и нестационарных условиях: непрерывно-временной графовой
динамики CT-TGNN/SDE-TGNN~(М1); федеративной оптимизации
FedLLM-API с дифференциальной приватностью~(М2);
многоуровневых темпоральных представлений TripleE-TGNN~(М3);
селективной модели пространства состояний MambaShield~(М4);
калиброванного по неопределённости вывода (стохастический
трансформер с UC-HGP, M5--M6); теоретико-игровой
сертификации FedGTD~(М7).

\item Разработать программное обеспечение, реализующее
предложенные модели и алгоритмы: сформулировать функциональные
требования, сценарии использования, архитектуру и провести
программную реализацию в виде единой платформы промышленного
уровня RobustIDPS.ai (DOI:~10.5281/zenodo.19129512;
17~нейросетевых моделей в 10~функциональных группах,
51~страница прогрессивного веб-приложения).

\item Провести экспериментальные исследования разработанных
моделей и алгоритмов на наборах данных-эталонах:
6~бенчмарках сетевого трафика общим объёмом свыше
84,2~миллиона размеченных образцов и 84~уникальных категорий
атак с использованием унифицированной 23-метрической системы
оценки (точность, калибровка, неопределённость, состязательная
устойчивость, вычислительная эффективность, приватность).

\item Интегрировать разработанное программное обеспечение в
реальный конвейер анализа сетевого трафика и защиты сетевой
инфраструктуры, подтвердив интеграцию актами внедрения от
организаций промышленного сектора.

\item Продемонстрировать, что разработанное программное
обеспечение допускает воспроизведение в другой практической
области, в качестве которой выбран сектор защиты бортовых
систем компьютерного зрения беспилотных летательных
аппаратов с трёхъярусной архитектурой
<<борт~--- дронопорт~--- облако>>.
\end{enumerate}

\subsection*{Методология и методы исследования}

В работе использовались методы системного анализа, теории
графов и спектральной теории графов, обыкновенных и
стохастических дифференциальных уравнений (СДУ),
функционального анализа, теории оптимального транспорта,
глубокого обучения, байесовского вывода, теории информации,
теории игр (равновесия Штакельберга и Нэша), теории
оптимизации, теории вероятностей, математической статистики,
теории дифференциальной приватности
(\emph{англ.}~Differential Privacy, ДП~--- формальная
криптографическая гарантия невозможности восстановления
индивидуальных записей из агрегированных результатов
обучения), а также численных и вычислительных экспериментов
на эталонных корпусах данных. Используемая методология
обеспечивает единое формальное описание гибридной СОПВ на
уровне кортежной модели и шести инвариантных свойств, что
позволяет доказательно увязать математические результаты
M1--M7 с прикладной задачей повышения устойчивости и
эффективности.

\subsection*{Научная новизна работы}

\begin{enumerate}
\item \textit{Единое математическое описание гибридной СОПВ.}
Впервые предложена кортежная формализация гибридной СОПВ
$\mathcal{H} = (\mathcal{V}_t, \mathcal{E}_t, X_t, A_t,
f_\theta, g_\psi, \pi)$ и согласованная с ней система из
шести инвариантных свойств П1--П6, охватывающая
состязательную устойчивость, робастность, эффективность,
распределённость, человеческий фактор и нестационарность.
На её основе построена матрица $3\times4$
<<условия$\times$требования>>, упорядочивающая семейство
M1--M7 как покрытие всех двенадцати ячеек.

\item \textit{Непрерывно-временная темпоральная графовая
динамика для гибридной СОПВ (CT-TGNN, SDE-TGNN).} Впервые в
составе гибридной СОПВ применена модель, в которой
представления узлов эволюционируют по нейронным обыкновенным
дифференциальным уравнениям на изменяющихся во времени
графах; сертифицированные границы устойчивости выводятся из
неравенства Грёнуолла с константой Липшица~$L_g$.
Стохастическое расширение SDE-TGNN заменяет
детерминированную динамику стохастическими дифференциальными
уравнениями Ито и аналитически распространяет первые два
момента (среднее и ковариацию) через уравнение
Фоккера--Планка, что даёт калиброванный сертифицированный
радиус через рандомизированное сглаживание классификатора.
Достигнуты: F1\,=\,99,0\,\%, ECE~=~0,042
(\emph{Expected Calibration Error}~--- ожидаемая ошибка
калибровки), сертифицированная устойчивость 76,2\,\% при
радиусе возмущения $R{=}0{,}25$ и кросс-доменный перенос
90,5\,\% F1 без дообучения.

\item \textit{Византийско-устойчивая федеративная оптимизация
для гибридной СОПВ (FedLLM-API,
\emph{Federated Large Language Model API}~--- федеративный
интерфейс программирования приложений больших языковых
моделей).} Разработан оригинальный федеративный оптимизатор,
совместно обеспечивающий покоординатное усечённое среднее,
формальную дифференциальную приватность с параметрами
$(\varepsilon{=}0{,}85,\ \delta{=}10^{-5})$ (где
$\varepsilon$~--- бюджет приватности, $\delta$~---
допустимая вероятность нарушения гарантии) и выравнивание
распределений гетерогенных клиентов через расстояние
Вассерштейна $W_2$ (метрика оптимального транспорта).
Доказана сходимость со скоростью
$\mathcal{O}(1/\sqrt{T})$ при 30\,\% повреждённых
участников; при 40\,\% византийских участников обеспечена
точность 87,1\,\% против 61,4\,\% у FedAvg
(\emph{Federated Averaging}~--- базовый алгоритм
федеративного усреднения). Интеграция большой языковой
модели (LLM, \emph{Large Language Model}) обеспечивает
87,6\,\% F1 при детектировании новых классов атак в режиме
\emph{zero-shot} (без дообучения).

\item \textit{Многоуровневые темпоральные графовые
представления для гибридной СОПВ (TripleE-TGNN, тройное
вложение).} Разработана архитектура тройных представлений
на уровнях сервисов, трассировок и узлов с двойным
темпорально-пространственным механизмом внимания и
\emph{упругой консолидацией весов} (EWC,
\emph{Elastic Weight Consolidation}~--- штраф на изменение
параметров, важных для ранее изученных задач), противодействующей
\emph{катастрофическому забыванию} при поступлении новых
данных. Достигнута точность 96,8\,\% при снижении
вычислительной стоимости на 40\,\% и сохранении 94,7\,\%
точности без дообучения при непрерывной эволюции топологии.

\item \textit{Селективные модели пространства состояний для
гибридной СОПВ (MambaShield).} Разработана архитектура с
быстрыми свёртками на основе \emph{быстрого преобразования
Фурье} (БПФ), снижающая сложность вывода с
$\mathcal{O}(L^2)$ до $\mathcal{O}(L\log L)$ (где $L$~---
длина последовательности) и интегрирующая прогрессивную
дистилляцию состязательной устойчивости. При 25\,\%
отравлении обучающих данных достигнута точность 91,4\,\%
против 73,8\,\% у незащищённой модели; PAC-байесовские
(\emph{Probably Approximately Correct, Bayesian}~---
вероятностные оценки обобщения, объединяющие частотную и
байесовскую парадигмы) сертификаты обобщения подтверждены с
погрешностью 2,7\,\%; задержка вывода 0,5~мс на один
сетевой поток.

\item \textit{Калиброванный по неопределённости вывод для
гибридной СОПВ (стохастический трансформер с UC-HGP,
\emph{Uncertainty-Calibrated Hierarchical Gaussian
Processes}~--- иерархические гауссовские процессы с
калиброванной оценкой неопределённости).} Впервые в составе
гибридной СОПВ разработана интеграция вариационного
байесовского механизма внимания с иерархическими
гауссовскими процессами, обеспечивающая декомпозицию
\emph{эпистемической} (модельной, сокращаемой при
дополнительных данных) и \emph{алеаторной} (неустранимой,
обусловленной шумом данных) неопределённости с
ECE\,=\,0,078 (снижение на 59\,\% относительно
детерминированного внимания) и сокращением нагрузки
аналитика SOC (\emph{Security Operations Centre}~--- центр
оперативного реагирования) на 43\,\% за счёт приоритизации
предупреждений по уровню неопределённости.

\item \textit{Теоретико-игровая сертификация гибридной СОПВ
(FedGTD, \emph{Federated Game-Theoretic Defence}~---
федеративная теоретико-игровая защита).} Построен
оригинальный механизм сертификации устойчивости через
\emph{равновесия Штакельберга} (стратегические равновесия
с лидером~--- защитником и ведомым~--- противником) с
алгоритмом мультипликативных весов, для которого доказана
\emph{граница сожаления}
$R(T)\le 2\sqrt{T\log|\mathcal{C}|}$ (где $T$~--- число
шагов, $|\mathcal{C}|$~--- мощность множества стратегий
защитника). Регуляризация согласованности связывает все семь
моделей с доказанной сходимостью. Точность 94,7\,\%
против адаптивных противников против 87,6\,\% у
нестратегических подходов.

\item \textit{Предметно-специфическая фундаментальная
модель (CyberSecLLM, \emph{Cybersecurity Large Language
Model}~--- большая языковая модель для кибербезопасности).}
Разработана архитектура с селективными блоками пространства
состояний, дополненными нейронными ОДУ, достигшая средней
zero-shot-точности 89,1\,\% на 11-задачном эталоне
CyberSecBench~--- улучшение на 20,5~процентных пункта
относительно GPT-4 при обработке последовательностей
графовых событий длиной в миллион токенов за 2,3~секунды.
\end{enumerate}

\subsection*{Теоретическая значимость}

Теоретическая значимость состоит в построении единого
математического описания гибридной СОПВ как объекта
исследования и интегрированной системы доказательств,
связывающей пять формальных дисциплин: нейронных
обыкновенных и стохастических дифференциальных уравнений
(CT-TGNN, SDE-TGNN, MambaShield); оптимального транспорта и
теории дифференциальной приватности (FedLLM-API);
селективных моделей пространства состояний и БПФ-свёрток
(MambaShield); вариационного байесовского вывода и
иерархических гауссовских процессов (стохастический
трансформер, UC-HGP); теории игр со стохастическим обучением
без сожаления (FedGTD). Совместное применение этих
аппаратов даёт формальные гарантии сходимости, устойчивости,
приватности и обобщения, верифицируемые как теоретически,
так и эмпирически, и впервые увязывает их в единый
научно-методический каркас повышения устойчивости и
эффективности гибридной СОПВ.

\subsection*{Практическая значимость}

Практическая значимость определяется разработкой моделей,
алгоритмов и специального программного обеспечения для
построения гибридной СОПВ, применимых в первую очередь в
секторе защиты сетевого трафика и в сектор защиты бортовых
систем беспилотных летательных аппаратов; кросс-доменно
проработана применимость в смежных секторах
(автономный транспорт, облачные и контейнерные среды).
Разработанная система обеспечивает:

\begin{itemize}
\item высокую точность обнаружения (94,7\,\% F1 на
CIC-IoT-2023) с одновременной состязательной устойчивостью
76,2\,\% при сертифицированном радиусе $R{=}0{,}25$;
\item сохранение приватности данных в распределённых
архитектурах за счёт федеративного обучения с
$(\varepsilon{=}0{,}85,\delta{=}10^{-5})$-дифференциальной
приватностью;
\item работу в режиме реального времени: пропускная
способность 8,7~млн событий/с при сквозной задержке 47~мс
и задержке вывода MambaShield 0,5~мс на поток;
\item совместимость с эволюционирующими протоколами, в том
числе с пост-квантовыми криптографическими стандартами
(модуль Multi-Agent PQC-IDS,
\emph{Post-Quantum Cryptography Intrusion Detection
System}~--- многоагентная система обнаружения вторжений в
постквантовом трафике, точность 96,2\,\%);
\item кросс-доменную переносимость 90,5\,\% F1 без
дообучения~--- подтверждено переносом из сектора анализа
сетевого трафика в сектор защиты бортовых систем БПЛА.
\end{itemize}

Практическая значимость подтверждена выпуском программной
платформы RobustIDPS.ai (DOI:~10.5281/zenodo.19129512),
интегрирующей \textbf{17~нейросетевых моделей} в
\textbf{10~функциональных групп} и \textbf{51~страницу}
прогрессивного веб-приложения (PWA,
\emph{Progressive Web Application}), четырёхуровневую
систему защиты больших языковых моделей (94,5\,\%
обнаружения по 517 сценариям атак с четырьмя коммерческими
бэкендами~--- Claude, GPT-4o, Gemini, DeepSeek), а также
модуль Multi-Agent PQC-IDS, и подтверждена двумя
полученными актами внедрения: от ООО~<<НИИ~ИКТ>>
(г.~Новосибирск; интеграция моделей обнаружения аномалий
сетевого трафика, устойчивых к атакам отравления, в
действующую систему обнаружения вторжений (СОВ) организации
и активное использование комплекса RobustIDPS.ai штатными
специалистами кибербезопасности) и от ООО~<<Ксайрикс>>
(г.~Москва; применение разработанных решений в составе
системы анализа сетевого трафика и защиты сетевой
инфраструктуры~--- обнаружение многостадийных атак и
координированных аномалий в графе сетевого взаимодействия,
выявление ранее неизвестных паттернов на основе анализа
неопределённости, адаптация к изменению статистических
характеристик трафика); акт от Центра искусственного
интеллекта НИЯУ МИФИ (г.~Москва; применение разработанного
каркаса AI-моделей гибридной СОПВ к защите бортовых систем
БПЛА) находится в стадии оформления.

\subsection*{Достоверность научных результатов}

Достоверность обеспечивается: строгой математической
формализацией гибридной СОПВ как кортежной модели и шести
свойств П1--П6 с доказательствами сходимости, устойчивости,
приватности и вычислительной сложности; экспериментальной
валидацией на шести открытых наборах данных общим объёмом
84,2~млн размеченных образцов и 84~уникальных категорий
атак; сравнительным анализом с современными базовыми и
передовыми методами с превосходством на 2,9--5,6 процентных
пункта; абляционными исследованиями, подтверждающими
измеримый вклад каждого компонента M1--M7; экспериментами
под шестью семействами состязательных атак (FGSM, PGD, C\&W,
DeepFool, гауссовский шум, маскирование меток); подтверждением
кросс-доменного переноса; представлением на международных
конференциях; публикациями в ведущих рецензируемых журналах
(Scopus / Web of Science); а также актами внедрения от двух
промышленных и одной исследовательской организации (третий
акт~--- в стадии оформления).

\subsection*{Апробация результатов исследования}

Основные положения и результаты были представлены и
обсуждены на следующих международных и российских научных
конференциях:

\begin{enumerate}
\item 2024 Conference of Young Researchers in Electrical and
Electronic Engineering (ElCon-2024), Санкт-Петербург, Россия, IEEE;
\item 2025 Conference of Young Researchers in Electrical and
Electronic Engineering (ElCon-2025), ETU-LETI, Санкт-Петербург,
Россия, IEEE;
\item International Conference ``Advances in Neural Computation,
Machine Learning, and Cognitive Research~IX''
(NEUROINFORMATICS-2025), Springer;
\item Третья Всероссийская научно-техническая конференция
<<Кибернетика и информационная безопасность>> (КИБ-2025),
НИЯУ МИФИ, Москва;
\item 2026 IEEE 27th International Conference of Young
Professionals in Electron Devices and Materials (EDM-2026),
Республика Алтай, Россия;
\item 11th IEEE Electronics System-Integration Technology
Conference (IEEE ESTC 2026), Хельсинки, Финляндия;
\item IEEE Research and Technologies for Society and Industry
(RTSI) 2026 Conference, Эспоо, Финляндия;
\item EMNLP 2026: The Annual Conference on Empirical Methods in
Natural Language Processing, Будапешт, Венгрия.
\end{enumerate}

\subsection*{Публикации}

Основные результаты опубликованы в 18~научных работах,
включая: 5~статей в рецензируемых журналах, индексируемых в
Scopus и Web of Science (Neurocomputing, Expert Systems with
Applications, Complex and Intelligent Systems, Optical Memory
and Neural Networks, IJASRE); 4~статьи в трудах международных
конференций (Scopus, IEEE~Xplore); 2~статьи, принятые к
публикации (Knowledge and Information Systems, Complex and
Intelligent Systems); 7~статей на рецензировании в ведущих
журналах IEEE~Transactions и других изданиях.

\subsection*{Основные научные результаты, выносимые на защиту}

Результаты, выносимые на защиту, перечислены в порядке,
отражающем их логические зависимости в рамках единого
математического описания гибридной СОПВ; каждый последующий
результат опирается на гарантии, установленные предыдущими.

\begin{enumerate}
\item Единое математическое описание гибридной СОПВ в виде
кортежа $\mathcal{H} = (\mathcal{V}_t, \mathcal{E}_t, X_t,
A_t, f_\theta, g_\psi, \pi)$ и системы шести инвариантных
свойств П1--П6, образующих матрицу $3\times4$
<<условия$\times$требования>>, на основе которой выстраивается
семейство M1--M7.

\item Модели CT-TGNN и SDE-TGNN сертифицированной
непрерывно-временной графовой динамики, в которых нейронные
ОДУ и стохастические дифференциальные уравнения управляют
эволюцией представлений узлов и рёбер на изменяющихся во
времени топологиях, а сертифицированные границы устойчивости
выводятся из неравенства Грёнуолла и аналитического
распространения моментов через уравнение Фоккера--Планка.

\item Алгоритм FedLLM-API византийско-устойчивой
федеративной оптимизации с доказанной сходимостью
$\mathcal{O}(1/\sqrt{T})$, формальной дифференциальной
приватностью $(\varepsilon{=}0{,}85,\delta{=}10^{-5})$ и
выравниванием распределений на основе оптимального
транспорта, обеспечивающий совместное обучение участников
гибридной СОПВ в условиях византийских нарушителей и
контролируемого бюджета приватности.

\item Архитектура TripleE-TGNN многоуровневых темпоральных
графовых представлений и архитектура MambaShield селективной
модели пространства состояний со сложностью
$\mathcal{O}(L\log L)$, совместно обеспечивающие выполнение
требований эффективности и устойчивости к отравлению
гибридной СОПВ при непрерывной эволюции топологии.

\item Модели стохастического трансформера и UC-HGP
калиброванного по неопределённости вывода с
эпистемически-алеаторной декомпозицией и сублинейной
границей сожаления для адаптации гибридной СОПВ к
нестационарным распределениям с поддержкой человеческого
фактора (триаж и приоритизация).

\item Теоретико-игровая система сертификации FedGTD,
интегрирующая гарантии результатов~(2)--(5) в единую
методологию оценки через равновесия Штакельберга и
регуляризацию согласованности, и предметно-специфическая
фундаментальная модель CyberSecLLM.

\item Программная реализация результатов в виде платформы
RobustIDPS.ai (17~нейросетевых моделей в 10~функциональных
группах и 51~страница веб-приложения), валидированная на
шести эталонных наборах данных (84,2~млн образцов,
84~категории атак) и подтверждённая актами внедрения
(2~получены, 1~в стадии оформления), а также инстанцирование
того же каркаса в смежном секторе защиты бортовых систем
БПЛА.
\end{enumerate}

\subsection*{Соответствие паспорту научной специальности}

Диссертация соответствует специальности 2.3.1 <<Системный
анализ, управление и обработка информации, статистика>>:

\noindent\textbf{П.4} <<Разработка методов и алгоритмов
решения задач системного анализа, оптимизации, управления,
принятия решений, обработки информации и искусственного
интеллекта.>>~--- соответствие подтверждается разработкой
семи оригинальных моделей и алгоритмов M1--M7 для задач
обнаружения и предотвращения вторжений и обработки
графовых потоков сетевого трафика в гибридной СОПВ.

\noindent\textbf{П.5} <<Разработка специального
математического и алгоритмического обеспечения систем
анализа, оптимизации, управления, принятия решений,
обработки информации и искусственного интеллекта.>>~---
соответствие подтверждается формализацией гибридной СОПВ в
виде кортежной модели и шести инвариантных свойств П1--П6,
а также реализацией единой программной платформы
RobustIDPS.ai с 17~интегрированными нейросетевыми моделями.

\subsection*{Структура и объём диссертации}

Диссертация состоит из введения, шести глав, заключения,
библиографического списка и трёх приложений. Основной текст
составляет 301~страницу. Библиографический список содержит
217~наименований. Глава~1 содержит системный анализ
гибридной СОПВ в прикладных секторах и формулирует пробелы.
Глава~2 даёт единое математическое описание гибридной СОПВ
и семейство моделей M1--M7. Глава~3 разрабатывает программную
платформу RobustIDPS.ai, включая технологический стек,
пакеты, версии и пользовательский интерфейс. Глава~4
представляет экспериментальную валидацию. Глава~5 описывает
внедрение и апробацию платформы в действующих организациях.
Глава~6 демонстрирует перенос разработанного каркаса в
сектор защиты бортовых систем БПЛА.

\section*{ОСНОВНОЕ СОДЕРЖАНИЕ РАБОТЫ}
\addcontentsline{toc}{section}{Основное содержание работы}

\textbf{Во введении} обоснована актуальность темы, определены
объект (гибридная СОПВ) и предмет (модели и алгоритмы
повышения устойчивости и эффективности), сформулирована
цель и восемь задач исследования, изложены научная новизна,
теоретическая и практическая значимость, основные положения,
выносимые на защиту, апробация результатов и структура
диссертации.

\textbf{Глава~1} проводит системный анализ гибридной СОПВ
в семи прикладных секторах и фиксирует пробелы, которые
невозможно закрыть в рамках одной модели. Анализ охватывает
три канонических типа обнаружения (сигнатурный,
поведенческий, протокольно-аналитический), выявляет их
ограничения, отдельно разбирает класс \emph{ИИ-усиленных}
коммерческих продуктов (CrowdStrike Charlotte~AI,
SentinelOne Purple~AI, Microsoft Security~Copilot,
Google Mandiant+Gemini, Cisco~Hypershield, IBM~QRadar~AI,
Palo~Alto Cortex~XSIAM) и фиксирует, что ни один из них
\emph{одновременно} не обеспечивает сертифицированной
устойчивости, приватной распределённости, эффективности и
калиброванной точности на динамических графовых структурах.
На рис.~\ref{fig:nidps_overview} приведён исполнительный
обзор сетевой инстанции гибридной СОПВ (NIDPS,
\emph{Network Intrusion Detection and Prevention System}~---
сетевая система обнаружения и предотвращения вторжений).

\begin{figure}[htbp]
\centering
\includegraphics[width=0.92\textwidth]{figures/fig_nidps_v1.png}
\caption{Исполнительный обзор сетевой инстанции гибридной
СОПВ (NIDPS): обучаемый контур поверх сигнатурных и
поведенческих детекторов.}
\label{fig:nidps_overview}
\end{figure}

\textbf{Глава~2} формализует гибридную СОПВ как кортеж
$\mathcal{H} = (\mathcal{V}_t, \mathcal{E}_t, X_t, A_t,
f_\theta, g_\psi, \pi)$ (см.~определения в разделе
``Актуальность'') и шесть инвариантных свойств
П1--П6 (рис.~\ref{fig:main_arch}). На этой основе строится
семейство M1--M7, образующее покрытие матрицы $3\times4$
<<условия$\times$требования>>.

\begin{figure}[htbp]
\centering
\includegraphics[width=0.96\textwidth]{figures/Main_arch_v1b.pdf}
\vspace{-0.5cm}
\caption{Архитектурная основа единого математического
описания гибридной СОПВ: семейство M1--M7 покрывает все
двенадцать ячеек матрицы <<условия$\times$требования>>.}
\label{fig:main_arch}
\end{figure}

\textit{Модель M1: CT-TGNN/SDE-TGNN~--- непрерывно-временная
графовая динамика для гибридной СОПВ.} Эволюция
представлений узлов задаётся параметризованным
дифференциальным оператором $f_\theta$ на изменяющейся во
времени топологии:
\begin{equation}
\frac{dh_v(t)}{dt} = f_\theta\bigl(h_v(t), A(t), X(t), t\bigr),
\end{equation}
где $h_v(t)$~--- скрытое состояние узла~$v$ в момент~$t$,
$A(t)$~--- матрица смежности, $X(t)$~--- матрица признаков
узлов. Сопряжённый метод чувствительности обеспечивает
обучение с константой $\mathcal{O}(1)$ по памяти;
\emph{темпоральная адаптивная пакетная нормализация}
согласует непрерывную динамику со слоями нормализации.

\textbf{Теорема~1 (Сходимость).} \emph{При константе Липшица
$L_f$ функции потерь и скорости обучения $\eta < 2/L_f$
процесс обучения сходится к стационарной точке со скоростью
$\mathcal{O}(1/\sqrt{T})$, где $T$~--- число итераций.}

\textbf{Теорема~2 (Сертифицированная устойчивость).}
\emph{Пусть векторное поле $f_\theta$ является
$L_g$-липшицевым по аргументу скрытого состояния. Тогда для
входных возмущений $\|\delta X\|\le\varepsilon$ выходная
чувствительность ограничена по неравенству Грёнуолла}
\begin{equation}
\|\delta h(T)\|\le\varepsilon L_g T\exp(L_g T).
\end{equation}

\begin{figure}[htbp]
\centering
\includegraphics[width=0.89\textwidth]{figures/fig_CT-TGNN_architecture.pdf}
\vspace{-0.5cm}
\caption{Архитектура CT-TGNN: непрерывно-временная графовая
динамика как обучаемый компонент гибридной СОПВ.}
\label{fig:CT-TGNN}
\end{figure}

\begin{figure}[H]
\centering
\includegraphics[width=0.96\textwidth]{figures/SDE-TGNN_arch.pdf}
\vspace{-1cm}
\caption{Архитектура SDE-TGNN с распространением неопределённости по Фоккеру--Планку}
\label{fig:SDE_TGNN}
\end{figure}

На \textbf{рис.~\ref{fig:SDE_TGNN}} показана архитектура SDE-TGNN с распространением неопределённости по Фоккеру--Планку: темпоральный граф проходит через кодировщик для формирования начального состояния $h_0$. Два параллельных пути: решатель СДУ вычисляет стохастические траектории на нескольких масштабах ($\tau_s{=}0.5$\,с, $\tau_m{=}2$\,с, $\tau_l{=}10$\,с), а модуль Фоккера--Планка аналитически распространяет среднее $\mu_t$ и ковариацию $\Sigma_t$ через связанные ОДУ. Представление, объединённое вниманием, поступает в классификатор, формирующий предсказания и калиброванную неопределённость. Сертифицированный радиус получен из вероятностей сглаженного классификатора.

Стохастическое расширение SDE-TGNN заменяет
детерминированную динамику стохастическими дифференциальными
уравнениями Ито и аналитически распространяет среднее и
ковариацию через уравнение Фоккера--Планка, давая
сертифицированный радиус через рандомизированное сглаживание
классификатора (рис.~\ref{fig:SDE_TGNN_robust}). Достигнуты:
F1\,=\,99,0\,\%, ECE\,=\,0,042, сертифицированная
состязательная устойчивость 76,2\,\% при $R{=}0{,}25$ и
кросс-доменный перенос 90,5\,\% F1.

\begin{figure}[htbp]
\centering
\includegraphics[width=0.80\textwidth]{figures/fig_sde_tgnn_robustness_vs_radius.pdf}
\vspace{-0.3cm}
\caption{SDE-TGNN: сертифицированная точность как функция
радиуса возмущения~$R$.}
\label{fig:SDE_TGNN_robust}
\end{figure}

\textit{Модель M2: FedLLM-API~--- федеративная оптимизация
гибридной СОПВ с приватностью.} Покоординатная агрегация
обновлений от $N$~участников с калиброванным гауссовским
шумом дисперсии~$\sigma^2$ и выравнивание распределений по
расстоянию Вассерштейна~$W_2$:
\begin{equation}
\begin{aligned}
\theta_{t+1}&=\text{TrimMean}\!\left(\{\theta_i^t\}_{i=1}^{N}\right)
+\mathcal{N}(0,\sigma^2 I),\\
W_2(\mu_i,\mu_j)&=\!\!\inf_{\pi\in\Pi(\mu_i,\mu_j)}\!\!
\sqrt{\int\!\!\|x-y\|^2 d\pi(x,y)},
\end{aligned}
\end{equation}
где $\theta_i^t$~--- локальное обновление параметров
участником~$i$ на шаге~$t$, $\mu_i$~--- эмпирическое
распределение признаков клиента~$i$, $\Pi(\mu_i,\mu_j)$~---
множество совместных распределений с маргиналами
$\mu_i$~и~$\mu_j$.

\textbf{Теорема~3 (Сходимость при повреждениях).}
\emph{Алгоритм сходится со скоростью
$\mathcal{O}(1/\sqrt{T})$ даже при 30\,\% повреждённых
участников при сохранении $(\varepsilon{=}0{,}85,
\delta{=}10^{-5})$-дифференциальной приватности.}

\begin{figure}[htbp]
\centering
\includegraphics[width=0.80\textwidth]{figures/fig_fedllm_api_architecture.pdf}
\vspace{-0.3cm}
\caption{Архитектура FedLLM-API: византийско-устойчивая
федеративная оптимизация гибридной СОПВ с дифференциальной
приватностью.}
\label{fig:FedLLM}
\end{figure}

Интеграция LLM даёт 87,6\,\% F1 при детектировании новых
аномальных категорий в режиме zero-shot; при 40\,\%
византийских участников точность 87,1\,\% против 61,4\,\%
у FedAvg.

\textit{Модель M3: TripleE-TGNN~--- многоуровневые
темпоральные графовые представления.} Двойное
темпорально-пространственное внимание (с весами
$\alpha_{uv}^{(t)}$ от узла~$u$ к узлу~$v$ в момент~$t$) и
упругая консолидация весов:
\begin{equation}
\alpha_{uv}^{(t)}=\frac{\exp(e_{uv}^{(t)})}
{\sum_{w}\exp(e_{vw}^{(t)})},\qquad
\mathcal{L}_{\text{EWC}}=\mathcal{L}_{\text{task}}
+\frac{\lambda}{2}\sum_i F_i(\theta_i-\theta_i^*)^2,
\end{equation}
где $e_{uv}^{(t)}$~--- логит внимания,
$\mathcal{L}_{\text{task}}$~--- потеря на текущей задаче,
$F_i$~--- диагональный элемент \emph{информации Фишера}
(оценка важности параметра~$\theta_i$ для ранее изученной
задачи), $\theta_i^*$~--- оптимальное значение параметра на
предыдущей задаче, $\lambda$~--- коэффициент регуляризации.

\begin{figure}[htbp]
\centering
\includegraphics[width=0.95\textwidth]{figures/fig_triplee_tgnn_architecture.pdf}
\vspace{-0.4cm}
\caption{Архитектура TripleE-TGNN: представления на уровнях
сервисов, трассировок и узлов с упругой консолидацией весов.}
\label{fig:TripleE_TGNN}
\end{figure}

Достигнута точность 96,8\,\% и снижение вычислительных
затрат на 40\,\% за счёт структурированной разреженности.

\textit{Модель M4: MambaShield~--- селективная модель
пространства состояний с защитой от отравления.} Дискретная
формулировка пространства состояний и БПФ-свёртки:
\begin{equation}
\begin{aligned}
h_{t+1}=\bar{A}h_t+\bar{B}x_t,\quad y_t=Ch_t+Dx_t,\\ \qquad
\mathcal{L}_{\text{distill}}=\text{KL}\!\left(p_{\text{st}}
\Big\|\frac{1}{K}\sum_k p_{\text{tch}_k}\right),
\end{aligned}
\end{equation}
где $h_t$~--- скрытое состояние, $x_t$~--- вход в момент~$t$,
$y_t$~--- выход; $\bar{A},\bar{B},C,D$~--- дискретизованные
матрицы пространства состояний; $\text{KL}(\cdot\|\cdot)$~---
расхождение Кульбака--Лейблера; $p_{\text{st}}$~---
распределение, предсказанное обучаемой моделью-учеником;
$p_{\text{tch}_k}$~--- распределение, предсказанное $k$-й
моделью-учителем из ансамбля размера~$K$.

\textbf{Теорема~4 (PAC-байесовское обобщение).} \emph{Для
обучающей выборки объёма~$N$, постериорного~$\rho$ и
приорного~$\pi$ распределений параметров истинная~$R$ и
эмпирическая~$\hat R$ потери удовлетворяют}
\begin{equation}
R(\theta)\le\hat{R}(\theta)
+\sqrt{\frac{\text{KL}(\rho\|\pi)+\ln(2N/\delta)}{2N}},
\end{equation}
\emph{с вероятностью не менее~$1-\delta$.}

\begin{figure}[htbp]
\centering
\includegraphics[width=0.65\textwidth]{figures/fig_mambashield_architecture.pdf}
\vspace{-0.3cm}
\caption{Архитектура MambaShield: $\mathcal{O}(L\log L)$
БПФ-свёртки с прогрессивной дистилляцией состязательной
устойчивости.}
\label{fig:mambashield}
\end{figure}

При 25\,\% отравлении достигнута точность 91,4\,\% против
73,8\,\% у незащищённой модели; задержка вывода 0,5~мс на
один сетевой поток.

\textit{Модели M5/M6: стохастический трансформер и
UC-HGP~--- калиброванный по неопределённости вывод.}
Вариационное байесовское внимание и иерархический
гауссовский процесс:
\begin{equation}
\begin{aligned}
q_\phi(\alpha)&=\mathcal{N}(\mu_\phi,\sigma_\phi^2),\qquad
\mathcal{L}=\mathbb{E}_{q_\phi}[\log p(y|x,\alpha)]
-\text{KL}(q_\phi\|p),\\
f(x)&\sim\mathcal{GP}\bigl(m(x),k(x,x')\bigr),
\end{aligned}
\end{equation}
где $q_\phi(\alpha)$~--- вариационное приближение
апостериорного распределения весов внимания~$\alpha$ с
параметрами~$\phi$ (среднее $\mu_\phi$, дисперсия
$\sigma_\phi^2$); $\mathcal{L}$~--- нижняя граница
маргинального правдоподобия (ELBO);
$\mathcal{GP}(m,k)$~--- гауссовский процесс с функцией
среднего~$m(\cdot)$ и ковариационным ядром~$k(\cdot,\cdot)$.

\begin{figure}[htbp]
\centering
\includegraphics[width=0.80\textwidth]{figures/fig_stochastic_transformer_architecture.pdf}
\vspace{-0.3cm}
\caption{Архитектура стохастического трансформера и UC-HGP:
эпистемически-алеаторная декомпозиция неопределённости и
триаж аналитика SOC.}
\label{fig:UCHGP}
\end{figure}

Достигнуты: ECE\,=\,0,078, кросс-доменный перенос 94,2\,\%,
сокращение нагрузки аналитика SOC на 43\,\%.

\textit{Модель M7: FedGTD~--- теоретико-игровая
сертификация гибридной СОПВ.} Равновесие Штакельберга с
алгоритмом мультипликативных весов:
\begin{equation}
w_i^{(t+1)}=w_i^{(t)}\exp(-\eta\ell_i^{(t)}),
\end{equation}
где $w_i^{(t)}$~--- вес стратегии~$i$ защитника на шаге~$t$,
$\ell_i^{(t)}$~--- наблюдаемая потеря этой стратегии,
$\eta$~--- скорость обучения.

\textbf{Теорема~5 (Граница сожаления).} \emph{Кумулятивное
сожаление защитника ограничено}
$R(T)\le 2\sqrt{T\log|\mathcal{C}|}$, \emph{где
$|\mathcal{C}|$~--- мощность множества стратегий защитника.}

\begin{figure}[htbp]
\centering
\includegraphics[width=0.80\textwidth]{figures/fig_game_theoretical_architecture.pdf}
\vspace{-0.3cm}
\caption{Архитектура FedGTD: теоретико-игровая сертификация
с регуляризацией согласованности для семейства M1--M7.}
\label{fig:FedGTD}
\end{figure}

Регуляризация согласованности объединяет все семь моделей
с доказанной скоростью сходимости и обеспечивает
94,7\,\% точности против адаптивных противников против
87,6\,\% у нестратегических подходов.

\textit{Фундаментальная модель CyberSecLLM.} ODE-дополненные
селективные блоки пространства состояний дают 89,1\,\%
среднюю zero-shot-точность на 11-задачном эталоне
CyberSecBench (улучшение на 20,5~процентных пункта над
GPT-4) при обработке последовательностей графовых событий
длиной в миллион токенов за 2,3~секунды.

\textbf{Глава~3} разрабатывает программную реализацию
гибридной СОПВ~--- платформу \textbf{RobustIDPS.ai}. В главе
сформулированы функциональные и нефункциональные требования,
сценарии использования, архитектура (FastAPI/PyTorch~2.x с
расширениями PyTorch~Geometric и DGL, React~18 / PWA,
PostgreSQL~16, Redis~7, Docker~Compose, плагин для
Suricata EVE-JSON), а также проведена программная
реализация. Сетевой трафик моделируется как темпоральный
граф $G_t=(\mathcal{V}_t,\mathcal{E}_t,X_t)$, где узлы~---
IP-адреса, хосты и сервисы, направленные рёбра~---
коммуникационные потоки, а матрица признаков
$X_t\in\mathbb{R}^{n_t\times d}$ кодирует 83~инженерных
признака пяти категорий. Унифицированная таксономия из
34~классов атак, формальный конвейер преобразования файлов
PCAP (\emph{Packet Capture}~--- сырой формат захвата
сетевых пакетов) и CSV-журналов в графовые структуры и
предметно-специфические конфигурации гиперпараметров
обеспечивают согласованную инстанциацию M1--M7.

\textbf{Глава~4} представляет комплексную экспериментальную
валидацию. Оценка охватывает шесть эталонных наборов
данных, объединённых в две группы общим объёмом
84,2~млн размеченных образцов и 84~категории атак
(табл.~\ref{tab:datasets}).

\begin{table}[H]
\centering
\caption{Эталонные наборы данных для экспериментальной
оценки.}
\label{tab:datasets}
\vspace{0.1cm}
\scriptsize
\begin{tabular}{@{}lccrcc@{}}
\toprule
\textbf{Набор данных} & \textbf{Домен} & \textbf{Год} & \textbf{Записи} & \textbf{Атаки} & \textbf{Призн.} \\
\midrule
\multicolumn{6}{@{}l}{\textit{Группа~1: общий сетевой трафик (IIS3D)}} \\[2pt]
CIC-IoT-2023 & IoT & 2023 & 46,6\,млн & 33 & 46 \\
CSE-CICIDS2018 & Корпорат. & 2018 & 16,2\,млн & 7 & 79 \\
UNSW-NB15 & Гибридный & 2015 & 2,5\,млн & 9 & 49 \\
\midrule
\multicolumn{6}{@{}l}{\textit{Группа~2: облако, микросервисы и периферия (ICS3D)}} \\[2pt]
Microsoft GUIDE & Облако & 2023 & 13,7\,млн & 15 & 51 \\
Container Security & Микросерв. & 2023 & 3,2\,млн & 8 & 93 \\
Edge-IIoT & Промышл. & 2022 & 2,0\,млн & 12 & 69 \\
\midrule
\textbf{Итого} & & & \textbf{84,2\,млн} & \textbf{84} & \\
\bottomrule
\end{tabular}
\end{table}

Используется унифицированная система из 23~метрик в шести
категориях: эффективность обнаружения (Accuracy, Precision,
Recall, F1, MCC, макро-F1); калибровка (ECE, MCE, оценка
Брайера, NLL); неопределённость (AURC, OOD~AUROC,
покрытие~/~риск); состязательная устойчивость (FGSM, PGD-20,
сертифицированная точность); вычислительная эффективность
(пропускная способность, задержка, число параметров);
приватность ($(\varepsilon,\delta)$-ДП, деградация
точности).

\begin{figure}[htbp]
\centering
\includegraphics[width=0.96\textwidth]{figures/Figure6_ROC_Generalization.pdf}
\caption{Сводная экспериментальная оценка гибридной СОПВ:
(а)~ROC-кривые на CIC-IoT-2023 (AUC\,0,991 против 0,974);
(б)~диаграммы надёжности по ECE; (в)~точность под
PGD-возмущениями; (г)~сертифицированная точность;
(д)~MambaShield под отравлением; (е)~тепловая карта
кросс-датасетной точности.}
\label{fig:roc_plot}
\end{figure}

\begin{table}[H]
\centering
\caption{Производительность единой гибридной СОПВ на
эталонных наборах данных.}
\small
\begin{tabular}{lcccc}
\toprule
\textbf{Набор данных} & \textbf{Точн.} & \textbf{F1} & \textbf{Базов.} & \textbf{Улучш.} \\
\midrule
CIC-IoT-2023 & 94,7\,\% & 94,7\,\% & 91,8\,\% & +2,9 \\
CSE-CICIDS2018 & 93,9\,\% & 93,9\,\% & 89,7\,\% & +4,2 \\
UNSW-NB15 & 93,8\,\% & 93,8\,\% & 88,2\,\% & +5,6 \\
MS\,GUIDE & 92,1\,\% & 92,0\,\% & 87,5\,\% & +4,6 \\
Container\,Sec & 91,5\,\% & 91,4\,\% & 86,9\,\% & +4,6 \\
Edge-IIoT & 89,8\,\% & 89,7\,\% & 84,3\,\% & +5,5 \\
\bottomrule
\end{tabular}
\end{table}

\textbf{Абляция:} удаление непрерывной динамики~--- $-6{,}8\,\%$;
многоуровневых представлений~--- $-8{,}3\,\%$; робастной
агрегации~--- $-21{,}4\,\%$; оптимального транспорта~---
$-15{,}8\,\%$; компонента LLM~--- $-53{,}4$~п.\,п.\ F1;
дистилляции~--- $-9{,}7\,\%$; вариационного внимания~---
$+114$ п.п.\ ECE (ухудшение калибровки).
\textbf{Состязательная устойчивость:} CT-TGNN сохраняет
76,2\,\% F1 при $\varepsilon{=}0{,}25$; единая система
достигает 89,6\,\% при комбинированных возмущениях.
\textbf{Кросс-доменный перенос:} 90,5\,\% F1 без дообучения
с сектора сетевого трафика в сектор защиты бортовых систем
БПЛА. \textbf{Производительность:} 8,7~млн событий/с,
задержка 47~мс. \textbf{Приватность:}
$(\varepsilon{=}0{,}85,\delta{=}10^{-5})$, деградация
точности лишь 3,1\,\%. \textbf{Неопределённость:}
ECE\,=\,0,078, сокращение объёма расследований на 43\,\%.
\textbf{Постквантовая совместимость:} модуль
Multi-Agent~PQC-IDS обеспечивает 96,2\,\% точности на
трафике постквантовых протоколов.

\textbf{Глава~5} описывает внедрение и апробацию платформы
\textbf{RobustIDPS.ai} (DOI:~10.5281/zenodo.19129512),
объединяющей \textbf{17~нейросетевых моделей} в
\textbf{10~функциональных групп} (AI~Command~Centre,
AI~Active~Defence, SOC~Intelligence, LLM~Attack~Surfaces,
MLSecOps~Standards, AI~Security \& Governance,
AI~Novel~Methods, AI~Data \& Models, Industry \& Research,
Notice~Board) и \textbf{51~страницу} прогрессивного
веб-приложения. Четырёхуровневая система защиты больших
языковых моделей достигает 94,5\,\% макро-усреднённого
обнаружения по 517 сценариям атак с четырьмя коммерческими
бэкендами (Claude, GPT-4o, Gemini, DeepSeek). Платформа
развёртывается как подключаемый модуль для существующих
сигнатурных СОВ (Snort, Suricata, Zeek).
Пользовательское исследование подтвердило сокращение времени
триажа аналитика на 34\,\%.

\begin{figure}[htbp]
\centering
\includegraphics[width=0.96\textwidth]{figures/Main_arch_v1a.pdf}
\vspace{-0.5cm}
\caption{Архитектура интегрированной платформы RobustIDPS.ai
с 17~моделями в 10~функциональных группах.}
\label{fig:main_arch_v1}
\end{figure}

\textbf{Внедрение:}

\begin{enumerate}
\item \textit{ООО~<<НИИ~ИКТ>>} (Научно-исследовательский
институт информационно-коммуникационных технологий,
г.~Новосибирск): интеграция в действующую систему
обнаружения вторжений (СОВ) организации моделей обнаружения
аномалий сетевого трафика, устойчивых к атакам отравления,
и активное применение программного комплекса RobustIDPS.ai
штатными специалистами кибербезопасности для
экспериментальных исследований протоколов информационной
безопасности и непрерывного анализа сетевого трафика
\textbf{(получен акт внедрения от 03.06.2026~г.).}

\item \textit{ООО~<<Ксайрикс>>} (г.~Москва): применение
разработанных решений в составе системы анализа сетевого
трафика и защиты сетевой инфраструктуры~--- обнаружение
многостадийных атак и координированных аномалий в графе
сетевого взаимодействия, выявление ранее неизвестных
паттернов аномального поведения на основе анализа
неопределённости, адаптация к изменению статистических
характеристик трафика \textbf{(получен акт внедрения,
2026~г.).}

\item \textit{Центр искусственного интеллекта НИЯУ МИФИ,
г.~Москва:} применение разработанного каркаса AI-моделей
гибридной СОПВ к защите бортовых систем БПЛА в рамках
трёхъярусной архитектуры <<борт~--- дронопорт~--- облако>>
\textbf{(акт внедрения находится в стадии оформления).}
\end{enumerate}

\textbf{Глава~6} переносит единый каркас гибридной СОПВ в
смежный сектор защиты бортовых систем беспилотных
летательных аппаратов с трёхъярусной архитектурой
<<борт~--- дронопорт~--- облако>>: бортовой ярус выполняет
инференс MambaShield с задержкой 0,5~мс на поток; дронопорт
агрегирует обновления через FedLLM-API; облачный ярус
реализует теоретико-игровую сертификацию FedGTD. Перенос
подтверждает доменную независимость разработанного
математического и алгоритмического каркаса.

\section*{ОСНОВНЫЕ РЕЗУЛЬТАТЫ И ВЫВОДЫ}
\addcontentsline{toc}{section}{Основные результаты и выводы}

Основные результаты диссертационного исследования
сформулированы в строгом соответствии восьми поставленным
задачам исследования.

\begin{enumerate}
\item Проведён систематический обзор подходов, применяемых в
современных гибридных СОПВ, и выявлены их недостатки и
уязвимости в прикладных секторах; задокументированы
критические отказы моделей и отставание устойчивости от
точности на 25--65~п.\,п.\ в условиях состязательных
воздействий.

\item Выполнен обзор методов повышения устойчивости и
обеспечения вычислительной эффективности; систематизированы
пять семейств моделей и алгоритмов; проведён сравнительный
анализ ИИ-усиленных коммерческих продуктов, который
зафиксировал отсутствие единой архитектуры, одновременно
обеспечивающей сертифицированную устойчивость, приватную
распределённость, эффективность и калиброванную точность
на динамических графовых структурах.

\item Сформулировано единое математическое описание
гибридной СОПВ в виде кортежа $\mathcal{H} = (\mathcal{V}_t,
\mathcal{E}_t, X_t, A_t, f_\theta, g_\psi, \pi)$ и системы
из шести инвариантных свойств П1--П6, образующих матрицу
$3\times4$ <<условия$\times$требования>>; в явном виде
поставлена многокритериальная задача повышения устойчивости,
точности, эффективности и приватности.

\item Разработаны семь моделей и алгоритмов M1--M7,
покрывающих все 12~ячеек матрицы:
непрерывно-временная графовая динамика
CT-TGNN/SDE-TGNN~(F1\,=\,99,0\,\%, ECE\,=\,0,042,
сертифицированная устойчивость 76,2\,\% при
$R{=}0{,}25$, кросс-доменный перенос 90,5\,\%);
федеративный алгоритм FedLLM-API с
$\mathcal{O}(1/\sqrt{T})$-сходимостью при 40\,\%
повреждённых участников
и $(\varepsilon{=}0{,}85,\delta{=}10^{-5})$-приватностью;
многоуровневые представления TripleE-TGNN (96,8\,\% F1,
снижение затрат на 40\,\%); MambaShield со сложностью
$\mathcal{O}(L\log L)$, достигающий 91,4\,\% F1 при
25\,\% отравлении; стохастический трансформер с UC-HGP
(ECE\,=\,0,078, $-43\,\%$ нагрузки SOC);
теоретико-игровая сертификация FedGTD (94,7\,\% F1
против адаптивных противников); доменно-специфическая
фундаментальная модель CyberSecLLM (89,1\,\%
zero-shot F1, $+20{,}5$~п.\,п.\ над GPT-4).

\item Разработано программное обеспечение~--- единая
платформа промышленного уровня RobustIDPS.ai
(DOI:~10.5281/zenodo.19129512) с
\textbf{17~нейросетевыми моделями} в
\textbf{10~функциональных группах} и
\textbf{51~страницей} прогрессивного веб-приложения,
четырёхуровневой защитой больших языковых моделей и
модулем Multi-Agent~PQC-IDS; 94,5\,\% обнаружения по
517 сценариям атак.

\item Выполнена экспериментальная валидация на
84,2~млн размеченных образцов и 84~категориях атак:
единая система достигает 94,7\,\% (CIC-IoT-2023),
93,9\,\% (CSE-CIC-IDS2018), 93,8\,\% (UNSW-NB15) и
89,8--92,1\,\% на остальных трёх наборах, превосходя
базовые методы на 2,9--5,6~п.\,п.

\item Подтверждена интеграция разработанного программного
обеспечения в реальный конвейер анализа сетевого трафика
и защиты сетевой инфраструктуры двумя полученными актами
внедрения: ООО~<<НИИ~ИКТ>> (г.~Новосибирск; интеграция
устойчивых к отравлению моделей в действующую СОВ и
активное использование RobustIDPS.ai штатными
специалистами кибербезопасности) и ООО~<<Ксайрикс>>
(г.~Москва; графовое обнаружение многостадийных атак,
оценка неопределённости, адаптация к дрейфу).

\item Подтверждена доменная независимость разработанного
каркаса переносом в сектор защиты бортовых систем БПЛА
(<<борт~--- дронопорт~--- облако>>); акт от Центра
искусственного интеллекта НИЯУ МИФИ, г.~Москва, находится
в стадии оформления и закрепит этот результат
документально.
\end{enumerate}

Совокупность результатов обеспечивает одновременное решение
задач состязательной устойчивости, робастности,
эффективности, распределённости, поддержки человеческого
фактора и нестационарности гибридной СОПВ через единое
математическое описание и согласованное семейство моделей и
алгоритмов M1--M7, объединяющее непрерывно-временную
графовую динамику, оптимальный транспорт, байесовский
вывод, селективные модели пространства состояний и теорию
игр.

\vspace{1cm}

\section*{ПУБЛИКАЦИИ ПО ТЕМЕ ДИССЕРТАЦИИ}
\addcontentsline{toc}{section}{Публикации по теме диссертации}

\sloppy

\subsection*{Статьи в журналах, индексируемых в Scopus / Web of Science}

\begin{enumerate}
\item \textit{Anaedevha~R.\,N., Trofimov~A.\,G., Borodachev~Y.\,V.}
Uncertainty-Calibrated Hierarchical Gaussian Processes for
Intrusion Detection with Multi-Scale Temporal Modeling //
Neurocomputing.~--- 2026.~--- V.\,677.~--- P.\,133105.~---
DOI:~10.1016/j.neucom.2026.133105 \textbf{(Q1)}.

\item \textit{Anaedevha~R.\,N., Trofimov~A.\,G., Borodachev~Y.\,V.}
MambaShield: Selective State-Space Models for
Poisoning-Resilient Sequential Modeling // Expert Systems
with Applications.~--- 2026.~--- P.\,131175.~---
DOI:~10.1016/j.eswa.2026.131175 \textbf{(Q1)}.

\item \textit{Anaedevha~R.\,N., Trofimov~A.\,G., Borodachev~Y.\,V.}
Temporal Adaptive Neural Ordinary Differential Equations
with Deep Spatio-Temporal Point Processes for Real-Time
Network Intrusion Detection // Complex and Intelligent
Systems.~--- Springer, 2026.~---
DOI:~10.1007/s40747-026-02291-7 \textbf{(Q1)}.

\item \textit{Anaedevha~R.\,N., Trofimov~A.\,G.} Improved
Robust Adversarial Model against Evasion Attacks on
Intrusion Detection Systems // Optical Memory and Neural
Networks.~--- 2024.~--- V.\,33(Suppl.\,3).~---
P.\,S414--S423.~--- DOI:~10.3103/S1060992X24700681
\textbf{(Q3)}.

\end{enumerate}

\subsection*{Труды конференций}

\begin{enumerate}
\setcounter{enumi}{5}
\item \textit{Anaedevha~R.\,N., Trofimov~A.\,G.} Integrated
Deep Q-Networks and Actor-Critic Models for Enhanced
Poisoning Attack Detection // 2025 Conf.\ Young Researchers
in Electrical and Electronic Engineering (ElCon).~---
IEEE~Xplore, 2025.

\item \textit{Anaedevha~R.\,N.} Application of Machine
Learning Models for Device Identification in Wireless
Network Traffic // 2024~ElCon.~--- IEEE~Xplore, 2024.~---
P.\,104--110.~--- DOI:~10.1109/ElCon61730.2024.10468413.

\item \textit{Anaedevha~R.\,N., Trofimov~A.\,G.} Stochastic
Multimodal Transformer with Uncertainty Quantification for
Robust Network Intrusion Detection // Advances in Neural
Computation, Machine Learning, and Cognitive Research~IX.
NEUROINFORMATICS~2025.~--- Springer.~---
DOI:~10.1007/978-3-032-07690-8\_35.

\item \textit{Anaedevha~R.\,N., Trofimov~A.\,G.}
SDE-TGNN: Multi-Scale Temporal Graph Neural Network with
Stochastic Dynamics for Cloud and Industrial Network
Intrusion Detection // Proc.\ 27th IEEE Int.\ Conf.\ Young
Professionals in Electron Devices and Materials
(EDM-2026), Altai Republic, Russia.~--- IEEE~Xplore, 2026.
\end{enumerate}

\subsection*{Принятые к публикации}

\begin{enumerate}
\setcounter{enumi}{9}
\item \textit{Anaedevha~R.\,N., Trofimov~A.\,G., Borodachev~Y.\,V.}
Stochastic PAC-Bayesian Transformers for Network Intrusion
Detection and NLP Applications // Knowledge and
Information Systems.~--- Springer, 2025.~--- Принята.

\item \textit{Anaedevha~R.\,N., Trofimov~A.\,G., Borodachev~Y.\,V.}
Byzantine-Resilient Stochastic Games for Federated
Multi-Cloud Intrusion Detection // Complex and Intelligent
Systems.~--- Springer, 2025.~--- Принята.
\end{enumerate}

\subsection*{На рецензировании}

\begin{enumerate}
\setcounter{enumi}{11}
\item \textit{Anaedevha~R.\,N., Trofimov~A.\,G.} TripleE-TGNN:
Triple-Embedding Temporal Graph Neural Networks for
Multi-Granularity Microservices Security // IEEE Trans.\
Dependable and Secure Computing.~--- 2026.

\item \textit{Anaedevha~R.\,N., Trofimov~A.\,G., Borodachev~Y.\,V.}
PQ-IDPS: Adversarially Robust Intrusion Detection for
Post-Quantum Encrypted Traffic // IEEE Trans.\ Neural
Networks and Learning Systems.~--- 2026.

\item \textit{Anaedevha~R.\,N., Trofimov~A.\,G., Borodachev~Y.\,V.}
FedLLM-API: Privacy-Preserving LLM Federation for Zero-Day
API Threat Detection // IEEE Trans.\ Information Forensics
and Security.~--- 2026.

\item \textit{Anaedevha~R.\,N., Trofimov~A.\,G., Borodachev~Y.\,V.}
CT-TGNN: Continuous-Time Temporal Graph Neural Networks
for Real-Time Threat Propagation // IEEE Trans.\ Services
Computing.~--- 2026.

\item \textit{Anaedevha~R.\,N., Trofimov~A.\,G., Borodachev~Y.\,V.}
Hybrid Spatial-Temporal Deep Learning for
Privacy-Preserving Encrypted Traffic Intrusion Detection //
IEEE Access.~--- 2026.

\item \textit{Anaedevha~R.\,N., Trofimov~A.\,G., Borodachev~Y.\,V.}
Differentially Private Optimal Transport for Multi-Cloud
Intrusion Detection // IEEE Access.~--- 2026.

\end{enumerate}

\end{document}
